\section{歴史的背景と問題意識 - 古典物理学の限界}

\noindent\textbf{この章の中心命題}
\[
\boxed{
u_{\mathrm{RJ}}(\nu,T)\xrightarrow[\nu\to\infty]{}\infty,\qquad
K_{\max}=h\nu-\Phi,\qquad
h\nu=E_n-E_m
}
\]
黒体放射、光電効果、原子スペクトルの三つの事実は、微視的世界では古典力学の状態概念と連続的なエネルギー交換の仮定をそのまま使えないことを示す。

19世紀末までに、ニュートン力学とマクスウェル電磁気学によって、巨視的な物理現象の大部分は説明されたと考えられていた。したがって、もし微視的世界も同じ枠組みで記述できるなら、粒子の状態は古典的な位相空間上の点として表され、時間発展も古典運動方程式で決まるはずである。しかし実際には、微視的現象ではこの期待が破れる。

\begin{experimentalfact}{黒体放射と紫外線発散}
\factitem{操作}{温度 $T$ の黒体から出る放射について、振動数 $\nu$ ごとのエネルギー密度 $u(\nu,T)$ を測定する。}
\factitem{前提}{黒体放射は、空洞内の電磁場が物質壁と熱平衡にある状態として扱える。}
\factitem{結果}{観測されるスペクトルは高振動数で減衰し、
\[
u_{\mathrm{P}}(\nu,T)=\frac{8\pi h\nu^3}{c^3}\frac{1}{e^{h\nu/(k_{\mathrm{B}}T)}-1}
\]
でよく記述される。}
\factitem{示唆}{高振動数の電磁場モードに等分配則をそのまま適用することはできない。}
\end{experimentalfact}

一方、古典論では、振動数 $\nu$ から $\nu+d\nu$ の間にある電磁場モード数と等分配則を用いると、単位体積・単位振動数あたりのエネルギー密度は
\[
u_{\mathrm{RJ}}(\nu,T)=\frac{8\pi \nu^2}{c^3}k_{\mathrm{B}}T
\]
となる。したがって全エネルギー密度は
\[
\int_0^\infty u_{\mathrm{RJ}}(\nu,T)\,d\nu = \infty
\]
と発散してしまう。この不一致が、エネルギー量子 $E=h\nu$ の導入を要請した。

\begin{experimentalfact}{光電効果}
\factitem{操作}{金属表面に振動数 $\nu$ の光を照射し、放出電子の有無と最大運動エネルギーを測定する。}
\factitem{結果}{電子が飛び出すためにはしきい振動数 $\nu_0$ が必要であり、放出電子の最大運動エネルギーは
\[
K_{\max}=h\nu-\Phi
\]
に従う。ここで $\Phi$ は仕事関数である。}
\factitem{示唆}{電子放出の有無を決める主要な量は、光の強度ではなく振動数である。}
\end{experimentalfact}

古典的な波動論だけで考えれば、電子へ渡されるエネルギーは光の強度に比例し、十分長く照射すればどの振動数でも電子は飛び出せるはずである。しかし上の結果では、電子放出の有無を決めるのは強度ではなく振動数であり、古典的波動像だけでは説明できない。

\begin{experimentalfact}{原子の安定性と線スペクトル}
\factitem{操作}{気体原子から放出または吸収される光を分光し、現れる振動数を測定する。}
\factitem{結果}{原子は安定に存在し、スペクトルは連続ではなく離散的な線として観測される。水素原子では線の波長が、リュードベリ定数 $R$ を用いて
\[
\frac{1}{\lambda}=R\left(\frac{1}{m^2}-\frac{1}{n^2}\right)
\qquad(n>m)
\]
の形に整理される。}
\factitem{示唆}{原子内部のエネルギー変化は連続量としてではなく、離散的な差として現れる。}
\end{experimentalfact}

古典電磁気学では、電荷をもつ粒子が加速するとラーモアの式
\[
P=\frac{e^2 a^2}{6\pi\varepsilon_0 c^3}
\]
に従って電磁波を放射する。したがって原子核のまわりを回る電子はエネルギーを失い、やがて原子核へ落ち込むはずである。しかし実際には原子は安定であり、しかもスペクトルは
\[
h\nu = E_n-E_m
\]
の形の離散線として理解される。ここでも古典理論の予測は観測と一致しない。

以上の事情から、古典力学における状態の記述および時間発展の法則をそのまま微視的世界へ持ち込むことはできない。したがって必要なのは、観測事実と両立するように状態・観測量・時間発展を新しく定める理論である。
